\subsubsection{Labels, routing domains, and preprocessing}
\label{ha:preprocessing}

Write $\mathcal F_i$ for the family of bags at scale $\delta_i$ and
$\mathcal B_i=\bigcup_{t\in T_i}\operatorname{Bag}_t$ for its balls.
The balls in $\mathcal B_i$ are pairwise disjoint, whereas balls at
different scales need not be disjoint.  Throughout this proof all
geometric sets are restricted to the continuous cube $Q=[0,1]^d$.
We abbreviate the common notation $B_Q$ to $B$ in this proof; in particular,
\[
 B(x,r)=\{y\in Q:\|x-y\|_1<r\},\qquad
 B(x,a\Vert b)=\{y\in Q:a<\|x-y\|_1<b\}
 \quad(0\le a<b).
\]
An ordinary ball includes its center; a hollow ball has strict inner
and outer boundaries.  Its width is $b-a$, and the width of $B(x,r)$
is $r$.  When a statement uses inner and outer radii for an ordinary
ball, they mean $0$ and $r$.  The sets $B(x,r)$ and
$B(x,0\Vert r)$ differ only at $x$; we keep this distinction when
testing membership.  In particular, every test below for an endpoint
in an input ball includes endpoints mapped to its center.
All centers used in the construction are centers of input balls, hence
belong to $\{0,1\}^d$: new hollow balls inherit an existing center.
Consequently, for each such center $x$, the map $y\mapsto\|x-y\|_1$
is affine on $Q$.  Intersections here mean intersections in the
continuous cube, not merely intersections on the finite set $X(V)$.

\begin{lemma}[Edge-disjoint paths from cut bounds]\label{ha:edge-paths}
If every cut separating distinct vertices $u,v$ of an undirected
multigraph has at least $q$ edges, where $q$ is a positive integer,
then the graph has $q$ pairwise edge-disjoint $u$--$v$ paths.
\end{lemma}
\begin{proof}
Give each edge an arbitrary orientation and maintain an integral
antisymmetric flow whose value on an oriented edge is in
$\{-1,0,1\}$.  Start with zero flow.  A directed traversal of an edge
is available when its signed flow in that direction is less than
$1$.  As long as the current flow value is less than $q$, augment
by one along an available directed $u$--$v$ path if one exists.
The edge bounds and flow conservation at all other vertices remain
valid, and the value increases by one.
If no such path exists, let $S$ be the vertices reachable from $u$
by available traversals.  Then $v\notin S$, and every edge leaving
$S$ has signed outward flow $1$.  Summing flow conservation over
$S$ says that the current flow value equals the number of those
edges, at least $q$, a contradiction.  Thus the procedure reaches
value $q$ after $q$ augmentations.
An integral flow of positive value contains a directed positive-flow
$u$--$v$ path: otherwise the vertices reachable on positive-flow
edges would have no positive outward flow, contradicting their
positive net flow.  Extract such a path and subtract one on its
edges.  Repeating extracts $q$ paths.  Each used edge had flow $1$
in its traversal direction and then has flow zero, so the extracted
paths are edge-disjoint.  Parallel edges are separate edges
throughout the argument.
\end{proof}

Let $\mathcal T[t]$ denote the set of hierarchy nodes in the subtree
rooted at $t$, including $t$.  Its leaves are identified with $V(t)$.
Recall that
\[
 P(t)=E(V(t),V\setminus V(t)),\qquad
 O(t)=E(V(t),V(t^*)\setminus V(t)).
\]
A ball record consists of its geometric ball and its labels; two
records may have the same geometric set.  Its owner $t(B)$ is a
hierarchy node.  A nonavoiding record will route in $G[V(t(B))]$.
An avoiding record also has a label $t_d(B)$, a strict descendant of
$t(B)$, and will route in
$G[V(t(B))\setminus V(t_d(B))]$.  In either case a further label
$t_P(B)$ is a set of strict descendants of $t(B)$.  Define the
\emph{conflict domain} and its leaf set by
\begin{align*}
 C(B)&=\mathcal T[t(B)]\setminus
       \bigcup_{s\in t_P(B)}\mathcal T[s]
       &&\text{if $B$ is nonavoiding},\\
 C(B)&=\mathcal T[t(B)]\setminus
       \left(\mathcal T[t_d(B)]\cup
       \bigcup_{s\in t_P(B)}\mathcal T[s]\right)
       &&\text{if $B$ is avoiding},\\
 L(B)&=C(B)\cap V.
\end{align*}
Thus $C(B)$ is a connected set of hierarchy nodes containing its root
$t(B)$.  The set $L(B)$, rather than the geometric points of $B$,
specifies the graph vertices available for charging.

\Needspace{13\baselineskip}
\begin{definition}[Valid labeling]\label{ha:valid-labeling}
A collection $Z$ of labeled ordinary and hollow balls is valid if the
following four conditions hold for every record $B$ with center $x$
and radii $a<b$.
\begin{enumerate}
 \renewcommand{\theenumi}{L\arabic{enumi}}
 \renewcommand{\labelenumi}{(\theenumi)}
 \item\label{ha:L1} The labels and conflict domain have the form just
 described; in particular $t(B)\in C(B)$.
 \item\label{ha:L2} For $s\in t_P(B)$ and $\{v,w\}\in P(s)$,
 \[
       \|x-X_v\|_1\ge b,\qquad \|x-X_w\|_1\ge b.
 \]
 \item\label{ha:L3} There is $u\in L(B)$ with $\|x-X_u\|_1\le a$
 and at least $k/4$ pairwise edge-disjoint paths starting at $u$ and
 ending at vertices of distance at least $b$ from $x$.  The paths lie
 in $G[V(t(B))]$ if $B$ is nonavoiding, and in
 $G[V(t(B))\setminus V(t_d(B))]$ if it is avoiding.
 \item\label{ha:L4} Distinct records whose geometric balls intersect
 have disjoint conflict domains.
\end{enumerate}
\end{definition}

The paths in \ref{ha:L3} may initially visit predecessor subtrees;
condition~\ref{ha:L2} will allow their relevant prefixes to avoid all
of them.  Also, the shape of the domains gives the useful equivalence
\begin{equation}\label{ha:domain-root-test}
 C(B)\cap C(B')\ne\varnothing
 \quad\Longleftrightarrow\quad
 t(B)\in C(B')\ \hbox{or}\ t(B')\in C(B).
\end{equation}
Indeed, a common node lies below both roots.  The roots are therefore
comparable, and the path from the higher root to that common node
contains the lower root and stays in the higher root's domain.
The converse follows since each root belongs to its own domain.

\paragraph{Predecessors and insertability.}
For $t\in T_i$ define
\[
 \operatorname{Pred}(t)=
 \{s:\ s\text{ is a strict descendant of }t,
                  \ s\in T_j\text{ for some }j<i\}.
\]
For $B=B(X_u,\delta_i)\in\operatorname{Bag}_t$, say that $B$ is
\emph{noninsertable by $s$} if $s\in\operatorname{Pred}(t)$ and an
endpoint $w$ of an edge of $P(s)$ satisfies
$\|X_u-X_w\|_1<\delta_i$.  A ball is \emph{insertable} if no
predecessor makes it noninsertable.  Every insertable input ball is
labeled nonavoiding, with
\[
       t(B)=t,\qquad t_P(B)=\operatorname{Pred}(t).
\]
Its center vertex $u$ belongs to $L(B)$.  To verify this, use the
assigned-bag property to choose an edge $\{u,v\}\in O(t)$ with
$v\notin V(t)$.  If $u\in V(s)$ for a predecessor $s$ of $t$, then
$\{u,v\}\in P(s)$.  Its endpoint $u$ is in the center-inclusive
ball $B$, contrary to insertability.  This argument also covers
distinct graph vertices having the same image under $X$.

\begin{lemma}[Cost of one predecessor]\label{ha:pruning}
Suppose $\delta_{i+1}\le\lambda\delta_i$ with $\lambda\le1/2$,
$|O(t)|\ge k\lambda|P(t)|$ for the selected hierarchy nodes, and each
bag is at least $C_3/k$-assigned.  For $s\in T_\ell$,
\[
 \sum_i\delta_i
   |\{B\in\mathcal B_i:B\text{ is noninsertable by }s\}|
 \le \frac{4}{C_3}\delta_\ell|\operatorname{Bag}_s|.
\]
\end{lemma}
\begin{proof}
The set being counted is empty for $i\le\ell$.  For each $i>\ell$,
assign to every counted ball one endpoint of an edge in $P(s)$ that
lies in that ball.  Pairwise disjointness of the center-inclusive
balls in $\mathcal B_i$ makes this assignment injective on endpoint
vertices.  There are at most $2|P(s)|$ such vertices, even when edges
are parallel or several vertices have the same image.  Thus the
left-hand side is at most
\begin{align*}
 2|P(s)|\sum_{i>\ell}\delta_i
 &\le \frac{2\lambda}{1-\lambda}|P(s)|\delta_\ell\\
 &\le 4\lambda|P(s)|\delta_\ell
 \le \frac{4}{k}|O(s)|\delta_\ell
 \le \frac{4}{C_3}|\operatorname{Bag}_s|\delta_\ell.
\end{align*}
The last inequality is precisely the assigned-bag cardinality bound.
\end{proof}

\Needspace*{3\baselineskip}
\begin{lemma}[Charging a valid labeled collection]\label{ha:labeled-energy}
Every finite collection $Z$ with a valid labeling satisfies
\[
       \frac{k}{4}\sum_{B\in Z}\operatorname{width}(B)
       \le \sum_{\{v,w\}\in E}\|X_v-X_w\|_1.
\]
\end{lemma}
\begin{proof}
Realize each graph edge $e=\{v,w\}$ by one fixed straight segment
from $X_v$ to $X_w$ in $Q$, of length $\ell_e=\|X_v-X_w\|_1$.
Different parallel edges are distinct resources even if their
segments coincide geometrically.  Concatenating these segments
realizes a graph path as a continuous curve.

Fix a record $B$ with center $x$ and radii $a<b$.  Take the paths in
\ref{ha:L3}, and stop each at its first vertex $z$ satisfying
$\|x-X_z\|_1\ge b$.  Every vertex in each resulting graph path
belongs to $L(B)$.  Avoidance of $V(t_d(B))$ already follows from
\ref{ha:L3}.  To check avoidance of $V(s)$ for $s\in t_P(B)$,
suppose the path first enters $V(s)$ along $\{v,w\}\in P(s)$.
The start vertex is outside $V(s)$, so this entering edge exists.
Its preceding endpoint $v$ occurs before the stopping vertex and
therefore has $\|x-X_v\|_1<b$.  This contradicts \ref{ha:L2}.
In particular, both graph endpoints of every edge of the stopped
path are in $L(B)$, including the edge containing its geometric
exit point.

On each of the resulting continuous curves, keep the portion between
the last point at radius $a$ preceding its first point at radius $b$
and that first point at radius $b$.  Those points exist by continuity.
The retained curve, apart from its endpoints, is in the open shell
$a<\|x-y\|_1<b$.  The function $y\mapsto\|x-y\|_1$ is
$1$-Lipschitz, so its length is at least $b-a$.  This also applies to
an ordinary ball by taking $a=0$; removing its center changes no
charged length.  Charge each retained portion to its underlying
graph edge.  The paths for one record are edge-disjoint, so this
record charges no portion of a graph edge twice.

Now fix a graph edge $e=\{v,w\}$.  If two distinct records charge
portions of $e$, the preceding paragraph gives
$v,w\in L(B)\cap L(B')$, so their conflict domains meet.
Condition~\ref{ha:L4} forces their geometric balls to be disjoint.
Hence their charged open portions of the fixed segment of $e$ are
disjoint.  Summing over all records charges $e$ at most $\ell_e$.
This reasoning uses edge portions, not the whole length of every
edge on a witness path.  It also handles zero-length edges, which
contribute no charge.  Each record receives at least
$(k/4)(b-a)$, and summing the edge budgets proves the claim.
\end{proof}

\begin{lemma}[Arrangement of insertable bags]\label{ha:insertable-arrangement}
Let $\mathcal T$ be a locally $k$-connected hierarchy, with selected
nonroot nodes $T$ satisfying $|O(t)|\ge k\lambda|P(t)|$ for $t\in T$.
Let $\mathcal F_i$ be a $\lambda$-geometric sequence of families
of $12C_3/k$-assigned bags of type $(\delta_i,T_i)$, where the
$T_i\subseteq T$ are pairwise disjoint.  Suppose all input balls are
insertable and
\[
 C_4\ge3,\qquad \lambda\le\frac{1}{6C_4},\qquad
 C_3\ge2\bigl((C_4+1)+4(C_4+2)^2\bigr).
\]
Then there is a finite valid collection $Z$ of labeled ordinary and
hollow balls, with centers inherited from the input balls, such that
\begin{equation}\label{ha:arrangement-width}
 \sum_{B\in Z}\operatorname{width}(B)
 \ge \frac{C_4}{12C_3}
       \sum_i\delta_i\sum_{t\in T_i}|\operatorname{Bag}_t|.
\end{equation}
\end{lemma}

The next parts prove Lemma~\ref{ha:insertable-arrangement}.  First we
show that it suffices for the assigned-bag proposition.

\begin{proof}[Reduction of Proposition~\ref{prop:hierarchy-assigned}
to Lemma~\ref{ha:insertable-arrangement}]
Let
$M=\sum_i\delta_i\sum_{t\in T_i}|\operatorname{Bag}_t|$
be the original total radius.  A ball can be noninsertable by several
predecessors; summing Lemma~\ref{ha:pruning} over all selected nodes
therefore bounds the total radius $N$ of noninsertable balls from
above, without needing a unique predecessor charge:
\begin{equation}\label{ha:preprocess-mass}
 N\le\frac{4}{C_3}M\le\frac14M.
\end{equation}
Here the original $24C_3/k$ assignment implies the weaker assignment
needed in that lemma, and the displayed condition on $C_3$ implies
$C_3\ge16$.

Delete all noninsertable balls.  If a bag retains at least half its
original balls, keep the remaining bag; otherwise delete the bag and
its node from the corresponding $T_i$.  A kept bag is
$12C_3/k$-assigned.  In a deleted bag more than half the original
balls were noninsertable, so its entire radius is at most twice its
noninsertable radius.  The same upper bound holds for the radius
removed from a kept bag.  The resulting sequence consequently has
total radius
\[
       M'\ge M-2N\ge M/2.
\]
Discard empty families, retaining the remaining scales in their
original order. The sequence retains disjointness at each scale
and the geometric scale condition. Deleting selected nodes can only remove predecessors,
so all surviving balls are still insertable after recomputing their
labels.  The hypotheses of Lemma~\ref{ha:insertable-arrangement}
therefore hold.  Combining that lemma with
Lemma~\ref{ha:labeled-energy} gives
\[
 \sum_{\{v,w\}\in E}\|X_v-X_w\|_1
 \ge\frac{k}{4}\frac{C_4}{12C_3}M'
 \ge\frac{k}{8}\frac{C_4}{12C_3}M,
\]
as required.
\end{proof}

\begin{lemma}[Validity of an insertable input ball]\label{ha:input-valid}
An insertable ball $B=B(X_u,\delta_i)\in\operatorname{Bag}_t$
with the input labels above satisfies \ref{ha:L1}--\ref{ha:L3}.
\end{lemma}
\begin{proof}
The definition gives \ref{ha:L1}, and insertability says exactly that
every endpoint in \ref{ha:L2} has distance at least $\delta_i$ from
$X_u$.  We have already proved that $u\in L(B)$.
Every selected node is nonroot, so $|O(t)|\ge k$ by the cut bound
in its parent's induced graph.  Moreover,
$|\operatorname{Bag}_t|\ge12C_3|O(t)|/k\ge12C_3>1$.
Choose another ball in the bag, with center vertex $w\in V(t)$.
Disjointness of the two open balls gives $X_w\notin B$, so
$\|X_u-X_w\|_1\ge\delta_i$.  Since $G[V(t)]$ is
$k$-edge-connected, Lemma~\ref{ha:edge-paths} supplies
$k$ pairwise edge-disjoint $u$--$w$ paths in this induced graph.
These witness \ref{ha:L3}, with inner radius zero, and in fact give
the stronger count $k$ in place of $k/4$.
\end{proof}
